\begin{afterword}
\summaryhead

The post asks the reader to imagine knowing nothing of biochemistry, so that the moving hand seems like magic. It quotes three passages by Lord Kelvin on the power of life to direct matter and on a ``vital principle.'' The post calls this vitalism: the view that an ``Élan vital'' made living matter move and took part in chemistry that non-living matter could not perform, until Wöhler's synthesis of urea dealt the theory a major blow. ``Élan vital'' was less an explanation than a curiosity-stopper. It gave the feeling of knowing, predicted nothing and fitted every outcome. The vitalists also revered their ignorance.

The author's principle is that ignorance exists in the map, not in the territory: mystery is a fact about a mind, not about a phenomenon. Vitalism, like phlogiston and Molière's ``dormitive potency,'' wrapped a mystery in a substance. The post calls such theories ``mysterious answers to mysterious questions'' and lists four signs: they stop curiosity, have no moving parts, are cherished by those who hold them, and leave the phenomenon as mysterious as before.

\responsehead

The post's principle is sound: mystery belongs to a state of knowledge, not to the thing studied. The author later credited nearly the same sentence to E. T. Jaynes. The problems are in the history the post offers as its illustration, and the post's own sources show them.

Its chief vitalist gave a chemical account of animal heat and motion. The ellipsis in the first passage cuts the sentence in which Kelvin compares the body to an electric motor run by a battery. The 1852 paper answers ``in the negative'' the question whether animals move matter by ``an inherent power,'' ascribes their heat and motion to the oxidation of food, and predicts that an animal climbing a hill would be ``sensibly less warm'' unless it breathed harder. Thompson's biography, the post's third source, says Kelvin's doctrine ``had little in common with the old crude vitalism which postulated a vital force.'' What Kelvin kept outside physics was the will's direction of those chemical forces and, in 1903, a ``Directive Power'' he tied to religion. That view is open to criticism of its own, and in the letters that followed in \textsc{The Times} Sir E. Ray Lankester denied that biologists were coming round to it. But it is not the theory the post describes.

The theory the post describes forbade something, which conflicts with the post's claim that it fitted every outcome. If vitalism held that some chemistry happens only in living matter, a laboratory synthesis could count against it, and it did not explain all outcomes equally. The story of that synthesis as ``a major blow'' is itself what the historian Douglas McKie called ``a chemical legend'': chemists went on asserting a vital force after 1828, and Wöhler's collaborator Liebig still doubted, decades later, that a leaf could grow by chemical forces.

The charge that ``the vitalists'' took ``pride in their ignorance'' is based on Kelvin alone, and the evidence that ``came knocking'' is not named. Five days later, in ``Failing to Learn from History,'' the author wrote of coming to realize ``how reasonable vitalism had seemed at the time.''

In short: a sound principle, that mystery is a fact about what someone knows, illustrated by a history of vitalism that the post's own sources do not support, since its chief vitalist explained animal heat by chemistry and the theory it describes forbade a result.
\end{afterword}
